\section{典范关系与自然边界条件}
\label{sec:canonical-relation-natural-boundary}

本节研究一般辛流形上的最小作用量原理，并在开弦和闭弦两种情境中，为多值作用量泛函的变分理论建立适当框架。

正如第 1 章中余切丛上的经典力学那样，考虑光滑道路集合
\[
 \mathcal PM=\{\gamma:[0,1]\to M\}.
\]
有自然的评价映射
$\operatorname{ev}_0,\operatorname{ev}_1:\mathcal PM\to M$，由
$\operatorname{ev}_i(\gamma)=\gamma(i)$（$i=0,1$）给出。把它们合并，定义映射
\[
 \operatorname{ev}=(\operatorname{ev}_0,\operatorname{ev}_1):
 \mathcal PM\longrightarrow M\times M.
\]

下面的典范关系概念在拉尔斯·赫尔曼德（Lars Hörmander）的傅里叶积分算子演算中曾起重要作用；参见 (Hor71) 和 (GS77)。

\begin{definition}[典范关系]\label{def:canonical-relation}
\stfhIndex{I073}{dian^fan^guan^xi}{典范关系（canonical relation）}
若子集 $C\subset M\times M$ 关于 $M\times M$ 上的辛形式
$-\omega\oplus\omega$ 是拉格朗日子流形，则称 $C$ 为\emph{典范关系}。
\end{definition}

下面将说明，典范关系在变分理论中扮演自然边界条件的角色。

在 $\mathcal PM$ 上定义\emph{作用 $1$-形式} $\alpha$：
\stfhIndex{I053}{zuo^yong^1^xing^shi}{作用 1-形式（action one-form）}
\begin{equation}
 \alpha(\gamma)(\xi)=\int_0^1\omega(\dot\gamma(t),\xi(t))\,dt.
 \label{eq:path-space-action-one-form}
\end{equation}
下面的公式揭示典范关系与最小作用量原理之间的联系。

\begin{proposition}\label{prop:action-one-form-differential}
设 $\alpha$ 是 $\mathcal PM$ 上的作用 $1$-形式，则
\begin{equation}
 d\alpha=\operatorname{ev}_1^*\omega-\operatorname{ev}_0^*\omega.
 \label{eq:action-one-form-boundary-differential}
\end{equation}
（这里把 $\mathcal PM$ 看作无限维流形。）
\end{proposition}

\begin{proof}
须证明对任意 $\gamma\in\mathcal PM$ 及
$\xi,\eta\in T_\gamma(\mathcal PM)$，
\begin{equation}
 d\alpha(\gamma)(\xi,\eta)
 =\omega(\xi(1),\eta(1))-\omega(\xi(0),\eta(0)).
 \label{eq:action-one-form-boundary-values}
\end{equation}
为计算 $d\alpha(\gamma)(\xi,\eta)$，选取二参数道路族
\[
 V:(-\epsilon,\epsilon)\times(-\epsilon,\epsilon)\longrightarrow\mathcal PM,
\]
满足
\[
 V(0,0)=\gamma,\qquad
 \left.\frac{\partial V}{\partial u}\right|_{(0,0)}=\xi,\qquad
 \left.\frac{\partial V}{\partial s}\right|_{(0,0)}=\eta.
\]
也以 $V=V(u,s,t)$ 表示相应映射
\[
 V:(-\epsilon,\epsilon)\times(-\epsilon,\epsilon)\times[0,1]
 \longrightarrow M.
\]
于是
\begin{equation}
 d\alpha(\gamma)(\xi,\eta)
 =\left.\left[
 \frac\partial{\partial u}\alpha\left(\frac{\partial V}{\partial s}\right)
 -\frac\partial{\partial s}\alpha\left(\frac{\partial V}{\partial u}\right)
 \right]\right|_{(0,0)}.
 \label{eq:path-space-exterior-derivative-definition}
\end{equation}
这里使用了恒等式
\[
 \left[\frac{\partial V}{\partial u},\frac{\partial V}{\partial s}\right]
 =V_*\left[\frac\partial{\partial s},\frac\partial{\partial u}\right]=0.
\]

\begin{remark}\label{rem:path-space-differential-forms}
严格说来，在道路空间 $\mathcal PM$ 上讨论微分形式及其外微分时需要更加谨慎。这里建议读者把公式\eqref{eq:path-space-exterior-derivative-definition}当作
$\mathcal PM$ 上 $2$-形式 $d\alpha$ 的定义。关于道路空间上微分形式的更系统讨论，参见 (GJP91) 和 (Chen73)。
\end{remark}

由作用 $1$-形式\eqref{eq:path-space-action-one-form}的定义，
\[
 \frac\partial{\partial u}\alpha\left(\frac{\partial V}{\partial s}\right)
 =\int_0^1\frac\partial{\partial u}
 \left(\omega\left(\frac{\partial V}{\partial t},
                         \frac{\partial V}{\partial s}\right)\right)dt.
\]
现固定任意\emph{无挠辛联络} $\nabla$，则
\[
\begin{aligned}
 \frac\partial{\partial u}
 \left(\omega\left(\frac{\partial V}{\partial t},
                         \frac{\partial V}{\partial s}\right)\right)
 &=\omega\left(\nabla_u\frac{\partial V}{\partial t},
                    \frac{\partial V}{\partial s}\right)\\
 &\quad+\omega\left(\frac{\partial V}{\partial t},
                    \nabla_u\frac{\partial V}{\partial s}\right).
\end{aligned}
\]
由 $\nabla$ 的无挠条件，
\[
 \nabla_u\frac{\partial V}{\partial s}
 =\nabla_s\frac{\partial V}{\partial u}.
\]
所以第二项会与\eqref{eq:path-space-exterior-derivative-definition}中求
$(\partial\alpha/\partial s)(\partial V/\partial u)$ 时出现的类似项抵消。于是只需计算
\[
\begin{aligned}
 \omega\left(\nabla_u\frac{\partial V}{\partial t},
                  \frac{\partial V}{\partial s}\right)
 &=\omega\left(\nabla_t\frac{\partial V}{\partial u},
                  \frac{\partial V}{\partial s}\right)\\
 &=\frac\partial{\partial t}
    \left(\omega\left(\frac{\partial V}{\partial u},
                          \frac{\partial V}{\partial s}\right)\right)
 -\omega\left(\frac{\partial V}{\partial u},
                  \nabla_t\frac{\partial V}{\partial s}\right).
\end{aligned}
\]
交换 $u,s$ 的角色并使用 $\nabla$ 的辛性质，得到
\[
\begin{aligned}
 &\frac\partial{\partial u}
   \left(\omega\left(\frac{\partial V}{\partial t},
                        \frac{\partial V}{\partial s}\right)\right)
 -\frac\partial{\partial s}
   \left(\omega\left(\frac{\partial V}{\partial t},
                        \frac{\partial V}{\partial u}\right)\right)\\
 &=2\frac\partial{\partial t}
   \left(\omega\left(\frac{\partial V}{\partial u},
                        \frac{\partial V}{\partial s}\right)\right)
 +\omega\left(\frac{\partial V}{\partial s},
                 \nabla_t\frac{\partial V}{\partial u}\right)
 -\omega\left(\frac{\partial V}{\partial u},
                 \nabla_t\frac{\partial V}{\partial s}\right)\\
 &=2\frac\partial{\partial t}
   \left(\omega\left(\frac{\partial V}{\partial u},
                        \frac{\partial V}{\partial s}\right)\right)
 -\frac\partial{\partial t}
   \left(\omega\left(\frac{\partial V}{\partial u},
                        \frac{\partial V}{\partial s}\right)\right)\\
 &=\frac\partial{\partial t}
   \left(\omega\left(\frac{\partial V}{\partial u},
                        \frac{\partial V}{\partial s}\right)\right).
\end{aligned}
\]
因此
\[
\begin{aligned}
 d\alpha(\gamma)(\xi,\eta)
 &=\int_0^1\left.
   \frac\partial{\partial t}
   \left(\omega\left(\frac{\partial V}{\partial u},
                        \frac{\partial V}{\partial s}\right)\right)
   \right|_{(0,0)}dt\\
 &=\int_0^1\frac\partial{\partial t}\omega(\xi(t),\eta(t))\,dt\\
 &=\omega(\xi(1),\eta(1))-\omega(\xi(0),\eta(0)),
\end{aligned}
\]
这正是\eqref{eq:action-one-form-boundary-values}。
\end{proof}

对给定子集 $C\subset M\times M$，记
\[
 \mathcal P_CM=\operatorname{ev}^{-1}(C)
 =\{\gamma\in\mathcal PM\mid(\gamma(0),\gamma(1))\in C\}.
\]
仍用同一字母 $\alpha$ 表示它在 $\mathcal P_CM$ 上的显然限制。由公式\eqref{eq:action-one-form-boundary-differential}立即得到如下推论，它把典范关系刻画为自然边界条件。

\begin{corollary}\label{cor:canonical-relation-natural-boundary}
作用 $1$-形式 $\alpha$ 在 $\mathcal P_CM$ 上闭，当且仅当
$C\subset M\times M$ 是典范关系。
\end{corollary}

这里给出几个重要的典范关系例子；它们会出现在后续关于弗洛尔理论的各章中。

\begin{example}\label{ex:canonical-relations}
\begin{enumerate}[label=(\roman*)]
\item \textbf{自由环路。}对角线 $\Delta\subset M\times M$ 是典范关系。相应道路空间 $\mathcal P_\Delta M$ 可自然等同于自由环路空间
\[
 \mathcal LM=\{\gamma:S^1\to M\mid S^1=\mathbb R/\mathbb Z\}.
\]
$\alpha$ 在 $\mathcal P_\Delta M$ 上的零点集是常值道路的集合，因而就是
$\mathcal P_\Delta M$ 中的一份 $M\hookrightarrow\mathcal P_\Delta M$。

\item \textbf{扭曲环路。}辛微分同胚 $\phi:M\to M$ 的图像
\[
 \Delta_\phi:=\{(x,\phi(x))\mid x\in M\}
\]
是典范关系。相应道路空间是
\[
 \mathcal L_\phi M
 =\{\gamma\in\mathcal PM\mid\phi(\gamma(0))=\gamma(1)\},
\]
其元素称为\emph{扭曲环路}。事实上，可以把一个 $\phi$-扭曲环路等同于映射圆
\[
 M_\phi:=[0,1]\times M/{\sim}\longrightarrow S^1
\]
的一个截面，其中等价关系 $\sim$ 由
\[
 (0,x)\sim(1,x')\quad\Longleftrightarrow\quad x'=\phi(x)
\]
给出。这时 $\alpha$ 的零点集无非是 $\operatorname{Fix}\phi$ 中的常值环路。

\item \textbf{拉格朗日边界条件。}考虑形如
\[
 C=L_0\times L_1
\]
的典范关系，其中 $L_0,L_1$ 是 $M$ 的拉格朗日子流形。$\alpha$ 的零点集对应于交集 $L_0\cap L_1$。
\end{enumerate}
\end{example}
